\section*{Introdução}

\addcontentsline{toc}{section}{Introdução}

\label{sec:introducao}

% (==>) Este é o fenômeno; esta é a questão; há evidência científica suficiente para justificar investigá-la

% ECONOMIA + INSTITUIÇÕES --> DECISÃO DE USO DA TERRA --> ESTOQUE AMBIENTAL
% Foco do Capítulo 1 (H1): como retornos econômicos esperados e o ambiente institucional condicionam a conversão da vegetação nativa.
% (A retroalimentação — vegetação × clima --> produtividade — constitui o objeto de H2 no Capítulo 2; este artigo investiga a formação e transformação do estoque).

% HIPÓTESE DA TESE INVESTIGADA NESTE ARTIGO (H1 e dimensões empíricas):
% ---
% H1: RETORNOS, INSTITUIÇÕES E CONVERSÃO
% Retorno econômico esperado + ambiente institucional --> decisão de uso da terra --> taxa de conversão de vegetação nativa
%
% DIMENSÃO EMPÍRICA: PREÇOS E CAPACIDADE PRODUTIVA LOCAL
% Choques favoráveis de preços agrícolas geram maior pressão de conversão onde a capacidade produtiva da terra é mais elevada.
%
% DIMENSÃO EMPÍRICA: CONDICIONAMENTO INSTITUCIONAL
% Regras, monitoramento, fiscalização e sanções condicionam a sensibilidade da taxa de conversão ao retorno econômico esperado.
%
% Heterogeneidade: respostas distintas entre atividades (lavoura vs. pastagem) e entre biomas (Amazônia, Cerrado e Pantanal).

% EIXOS CONCEITUAIS QUE CONECTAM O ARTIGO À TESE:
% ---
% ENDOGENEIDADE AMBIENTAL
% O meio ambiente é tratado como componente endógeno do processo econômico: seu estoque é transformado pelas decisões produtivas sob incentivos e instituições.
%
% GOVERNANÇA TERRITORIAL
% Instituições, regulação e mecanismos de implementação operam sobre agentes e territórios heterogêneos, condicionando a intensidade da resposta aos incentivos de mercado.
%
% INTERTEMPORALIDADE
% Decisões presentes de conversão de uso da terra alteram estoques ambientais acumulados e condicionam oportunidades e riscos econômicos futuros.

% ------------------------------------
%%% (I) CONTEXTO
% ------------------------------------
% Mapeamento bibliográfico de H1:
% - Dilema conservação vs. rentabilidade e determinantes da decisão privada: Angelsen (2010), Lubowski et al. (2008), Jung et al. (2022).
% - Implementação regulatória e enforcement local: Viallon et al. (2019), Assunção et al. (2023), Pacheco & Meyer (2022).
% - Preços relativos, custos e contexto territorial: Harding et al. (2021), Chomitz & Thomas (2003).
% - Reorganização produtiva vs. expansão de fronteira: Bragança (2018).
% - Mato Grosso como contexto empírico e diversidade de trajetórias: VanWey et al. (2013), Simões et al. (2020).

%%% REF1
%% Conflito: preservar ou converter (dilema econômico e determinantes privados/institucionais da decisão).
A expansão agropecuária costuma ser analisada à luz do dilema entre conservar a vegetação nativa, promovendo sustentabilidade ambiental, ou converter a terra para explorar usos economicamente mais rentáveis \cite{Angelsen2010}. Essa representação, contudo, simplifica uma decisão que não depende apenas da aptidão agronômica da terra, da expectativa de sua valorização ou da rentabilidade econômica esperada, mas também é influenciada pelo ambiente institucional no qual ela se encontra inserida \cite{LubowskiEtAl2008}. A decisão privada de converter vegetação é condicionada por preços, custos de produção, restrições de crédito, expectativas, liquidez, direitos de propriedade, regras ambientais, \textit{enforcement} e probabilidade de sanções pelo uso irregular da terra \cite{JungEtAl2022}. Nesse ambiente, a definição e a proteção dos direitos de propriedade influenciam os incentivos dos agentes para o uso da terra \cite{alstonTitlesConflictLand1999}. 

%%% REF2
%% Ponte para H1: retornos econômicos, incentivos institucionais e evolução do estoque de vegetação nativa.
A par disso, aumentos nos preços dos produtos agropecuários podem elevar a pressão sobre as florestas, enquanto políticas de controle do desmatamento podem modificar essa resposta, com efeitos distintos a depender do instrumento, dos usos da terra e das áreas afetadas \cite{HardingEtAl2021}. Sob essa perspectiva, o monitoramento e a fiscalização podem conter o desmatamento, como indica a evidência para a Amazônia brasileira \cite{AssuncaoEtAl2023}. Considerados em conjunto, esses argumentos permitem tratar o estoque de vegetação nativa como um componente ambiental endógeno ao processo econômico: sua conservação ou conversão depende, em parte, das decisões de uso da terra tomadas pelos agentes diante dos retornos esperados, dos direitos de propriedade e dos custos e riscos associados às regras ambientais e à sua implementação. A evolução desse estoque constitui, assim, parte do próprio processo de alocação econômica da terra, condicionado pelo ambiente institucional. %% Heterogeneidade institucional e implementação local (regras formais vs efetividade da fiscalização/enforcement).

%%% REF3
%% Justificativa para H1
Importa considerar, ainda, que a resposta a esses incentivos pode envolver tanto a conversão de vegetação nativa quanto a reorganização da produção em áreas já convertidas, como ocorreu na bacia do Tapajós\footnote{Abrange principalmente os estados brasileiros de Mato Grosso e do Pará, além de pequenas porções no Amazonas e em Rondônia.}, onde aumentos nos preços relativos \textit{crop-to-beef} estiveram associados à substituição de pastagens por lavouras e a menor desmatamento local \cite{Braganca2018}. Na Amazônia, as características biofísicas e o acesso aos mercados estão associados a diferenças no uso agropecuário da terra \cite{ChomitzThomas2003}. \citeonline{BermanEtAl2023} evidenciam que, nos trópicos, a resposta do desmatamento a alterações nos preços agrícolas também varia conforme o contexto econômico e territorial sobre o qual os incentivos incidem. É necessário, portanto, investigar quando e com que intensidade variações nos preços dos produtos agropecuários se traduzem em conversão da vegetação nativa e como essa relação é condicionada pela capacidade produtiva local da terra e pelo ambiente institucional.

%%% REF4
%% Mato Grosso como caso empírico: grande escala agropecuária, três biomas e diferenciais biofísicos e econômicos.
Nesse contexto, Mato Grosso constitui um caso empírico especialmente informativo para investigar como os incentivos econômicos para usos agropecuários e o ambiente institucional se relacionam com a conversão da vegetação nativa, o que possibilita uma análise detalhada das dinâmicas de uso da terra em um ambiente heterogêneo. O estado combina produção agropecuária diversificada em larga escala com a presença dos biomas Amazônia, Cerrado e Pantanal \cite{VanWeyEtAl2013}, e a literatura documenta, adicionalmente, tanto a expansão e a intensificação agrícola quanto a diversidade das trajetórias de uso da terra em Mato Grosso \cite{SimoesEtAl2020}.

% ------------------------------------
%%% (II) PROBLEMA, OBJETIVOS E JUSTIFICATIVA
% ------------------------------------

%%% REF5
%% Problema central e formulação da questão de pesquisa (H1 em Mato Grosso).
Diante desse quadro, o problema central deste artigo consiste em compreender como os retornos econômicos esperados e o ambiente institucional condicionam a conversão da vegetação nativa em territórios ambiental e economicamente heterogêneos, utilizando Mato Grosso como caso de estudo. O \textbf{problema de pesquisa} pode ser sintetizado no seguinte questionamento:

\medskip

\noindent\emph{Como os retornos econômicos esperados e o ambiente institucional condicionam a taxa de conversão da vegetação nativa em Mato Grosso, e em que medida essa relação varia conforme a capacidade produtiva local da terra, as atividades agropecuárias e os biomas?}

\medskip

%%% REF6
%% Desdobramento do problema nas três dimensões empíricas relacionadas a H1: (i) preços × capacidade produtiva; (ii) condicionamento institucional; (iii) heterogeneidade entre atividades e biomas.
Essa pergunta desdobra-se em três dimensões empiricamente relacionadas: (i) se aumentos nos preços dos produtos agropecuários se traduzem em maior conversão onde a capacidade produtiva da terra para esses usos é mais elevada; (ii) se as condições institucionais modificam a intensidade dessa resposta; e (iii) se essas relações entre expectativas e respostas aos ambientes econômico e institucional variam entre atividades agropecuárias e biomas. 

%% Objetivo geral e objetivos específicos (i)–(iv): conversão bruta, resposta a preços × capacidade, instituições e heterogeneidade.
O \textbf{objetivo geral} deste artigo é investigar como e com que intensidade os determinantes observáveis do retorno econômico esperado e o ambiente institucional se relacionam com as taxas de conversão da vegetação nativa em Mato Grosso. Os \textbf{objetivos específicos} são: (i) construir medidas de conversão bruta da vegetação nativa para agricultura e pastagem; (ii) examinar a resposta das taxas de conversão a variações nos preços dos produtos agropecuários em interação com a capacidade produtiva local predeterminada; (iii) parametrizar e avaliar o condicionamento institucional dessa resposta, explicitando os limites de interpretação causal; e (iv) verificar a heterogeneidade dessas respostas entre atividades agropecuárias e biomas.

% ------------------------------------
%%% (III) HIPÓTESE
% ------------------------------------

%%% REF7
%% H1 — Governança territorial e heterogeneidade espacial: investiga como incentivos econômicos e instituições condicionam as decisões de conversão ou preservação da vegetação, considerando diferenças entre territórios e biomas.

%% Enunciação da hipótese H1 da tese, examinada por dimensões empíricas articuladas.
A partir dessa delimitação, o artigo examina a seguinte hipótese:

% -------------------------------------------------------------------------
\medskip

% Hipótese geral do artigo (H1):
% Maiores retornos econômicos esperados dos usos agropecuários elevam a taxa de conversão de vegetação nativa.
\noindent\textbf{(H1) Retornos, instituições e conversão.}

\noindent\emph{Elevações nos retornos econômicos esperados dos usos agropecuários aumentam a taxa de conversão da vegetação nativa, sendo essa resposta, \textit{ceteris paribus}, amplificada pela capacidade produtiva local da terra e modulada pela efetividade do ambiente institucional.}

\medskip

%% DESENHO PRELIMINAR H1
% Conversão da vegetação nativa = f(preços agropecuários, produtividade da terra, ambiente institucional).
% Recorte candidato: municípios de Mato Grosso, 2001--2024, conforme a cobertura dos dados.
% Os arquivos abaixo estão em 2_INTI/xdata/chap1/processed/tabular/.
%
% [Y] conversão da vegetação nativa: área convertida para agricultura ou pastagem / área nativa no início do período, por destino.
% - Fonte: MapBiomas, Coleção 10.1.
% - Arquivo disponível: h1_h2__mapbiomas_classes__col10_1__municipio_bioma_anual.csv.
% - Pendente: construir as transições anuais; o arquivo contém áreas por classe, não conversões.
%
% [X1] preços agropecuários: preços internacionais de soja e milho (USD/t), disponíveis mensalmente.
% - Fonte: Banco Mundial, Pink Sheet.
% - Arquivo: h1__prices__worldbank__mensal_robustez.csv (preco_nominal; price_soja e price_milho).
%
% [X2] produtividade da terra: rendimento observado de soja e milho, por município e ano (kg/ha).
% - Fonte: IBGE, Produção Agrícola Municipal (PAM).
% - Arquivo: h1_h2__pam_soja_milho__ibge__municipio_anual.csv (rendimento_medio_kg_ha).


% - Alternativa a avaliar: rendimento potencial de soja e milho, estimado pelo GAEZ/FAO-IIASA.
% - Arquivo: h1__suitability_yield__gaez_v4_6190__candidatos.csv (variantes yx, em kg de matéria seca/ha).
% - PAM mede rendimento nas áreas colhidas; GAEZ representa potencial sob cenários definidos.
% - xx_rendsoja-rendmilho = rendimento potencial de soja e milho, kg-MS/ha, municipal, 1961-1990
% 2_INTI\xdata\chap1\processed\tabular\h1__suitability_yield__gaez_v4_6190__candidatos.csv
% * MS: Matéria seca; 
% ** Observar os códigos: (I) sx...: índice de aptidão adimensional, na escala 0–10.000; 
% (II) yx...: kg de matéria seca por hectare (kg MS/ha), medida de rendimento potencial.

%
% [X3] ambiente institucional: número de autos de infração do IBAMA, por município e ano.
% - Fonte: IBAMA, registros de autuação.
% - Arquivo: h1__autos_multas_ibama__municipio_anual_candidato.csv (n_autos_ibama_it).
% - Pendente: validar as infrações pertinentes e a cobertura; a contagem não mede diretamente a efetividade da fiscalização.
% As autuações também podem responder ao desmatamento, o que limita sua interpretação causal.
% Outras fontes institucionais dependerão das lacunas identificadas.
%
% Primeiro: tratar os dados, definir as medidas e realizar a análise descritiva.
% Depois: avaliar modelos de efeitos fixos, defasagens, logs e interações, conforme os dados e as previsões de H1.
% A relação acima organiza as variáveis candidatas; a especificação econométrica será definida nessa etapa posterior.

%% Mapeamento manual de referência, anterior à consolidação de H1:
% 2_INTI\chapters\chap1\stats\HIPOTESES TESE - CHAP1.txt

%% Delimitação conceitual e variável dependente: vegetação nativa, conversão agropecuária e taxa de conversão.
Nessa formulação, a vegetação nativa corresponde à cobertura vegetal natural existente no território, enquanto sua conversão designa a substituição dessa cobertura por usos agropecuários, como lavouras ou pastagens. A taxa de conversão, por sua vez, refere-se à parcela da vegetação nativa disponível no início de determinado período que é convertida para esses usos no mesmo período, constituindo o principal resultado de interesse. A ausência de conversão observada, contudo, não equivale necessariamente a uma decisão ativa de conservação ou à integridade ecológica da vegetação remanescente. 

Por seu turno, o retorno econômico esperado dos usos agropecuários corresponde ao benefício privado que se espera obter com a utilização produtiva da terra, considerando as receitas, os custos e a produtividade associados à atividade. Sua relevância para a decisão de converter depende também da atratividade dos usos alternativos da terra, dos custos da conversão e das restrições institucionais. O exame empírico de H1 articulará os preços agropecuários e medidas predeterminadas de capacidade produtiva da terra às condições institucionais locais.

%%% REF8
%% Dimensão empírica de H1: resposta aos preços conforme a capacidade produtiva local.
Quanto à capacidade produtiva local, pressupõe-se que aumentos de preços produzam maior pressão sobre a conversão onde a terra oferece maior aptidão para os usos beneficiados, conforme discutido na literatura apresentada a seguir. A investigação empírica examinará uma implicação observável de H1 por meio da interação entre preços agropecuários e medidas de produtividade da terra, sem estimar diretamente o retorno líquido esperado, uma vez que fatores não observados, como restrições de crédito, liquidez e aversão ao risco, também influenciam a decisão do produtor.

%%% REF9
%% Dimensão empírica de H1: condicionamento institucional da resposta aos incentivos econômicos.
A capacidade produtiva, entretanto, não é o único fator que modula essa resposta: as regras de uso da terra e sua implementação alteram os custos e os riscos associados à abertura de novas áreas. O ambiente institucional compreende as regras formais, informais e os mecanismos de implementação que modificam os incentivos econômicos à conversão, incluindo direitos de propriedade, regulação ambiental, monitoramento, fiscalização e sanções. Não se pressupõe, ressalte-se, que essas instituições determinem isoladamente a conversão, mas que condicionem a intensidade com que alterações nos retornos econômicos esperados se traduzem em conversão efetiva. Quando uma medida representar, de maneira defensável, maior efetividade da fiscalização, espera-se menor sensibilidade da conversão ao retorno, mantidas as demais condições relevantes. 

%%% REF10
%% Heterogeneidade empírica (atividades e biomas), escopo da análise e critérios para resultados nulos ou contrários.
Consoante o contextualizado anteriormente, as relações consideradas no exame de H1 são comparadas com vistas a verificar a existência de heterogeneidade entre atividades agropecuárias e biomas, porquanto diferentes atividades apresentam tecnologias, estruturas de custos e retornos distintos, enquanto os biomas diferem em características biofísicas e produtivas. Essa heterogeneidade é tratada como dimensão empírica da análise, sem constituir uma hipótese autônoma, sendo interpretada à luz das diferenças de retorno econômico esperado, capacidade produtiva local e ambiente institucional.

%%% REF11
%% Justificativa artigo em relação à literatura e pesquisa empírica existentes: contribuições pretendidas.
Em termos empíricos, o artigo se justifica porque a literatura aponta respostas distintas do uso da terra aos incentivos de mercado, ora associando altas de preços à maior pressão sobre florestas, ora indicando substituição de pastagens por lavouras ou deslocamento espacial do desmatamento conforme os instrumentos de conservação vigentes \cite{Braganca2018,HardingEtAl2021,BermanEtAl2023}. Como será aprofundado na seção seguinte, essas divergências exigem distinguir elevações no nível de preços agrícolas de mudanças nos preços relativos entre atividades, bem como separar a abertura de novas áreas (conversão bruta) da reorganização produtiva em terras já abertas. Esses resultados tornam pertinente investigar, em Mato Grosso, em quais contextos territoriais e institucionais os incentivos produtivos efetivamente resultam em conversão de vegetação nativa.

%% Relevância teórica e eixos da tese: meio ambiente endógeno, dimensão territorial e intertemporal.
Em termos conceituais, o trabalho trata o estoque de vegetação nativa como um componente essencial \textbf{endógeno} ao processo econômico, cuja evolução depende, em parte, de decisões de conversão ou conservação da terra condicionadas pelos retornos econômicos esperados e pelo ambiente institucional. Por sua vez, o estoque remanescente pode modificar as condições produtivas e os retornos que orientam decisões posteriores. A dimensão institucional da \textbf{governança territorial} é examinada ao investigar como as regras de uso da terra e sua implementação influenciam a resposta da conversão da vegetação aos retornos econômicos esperados. Essa abordagem considera que uma mesma regra formal pode produzir respostas distintas a depender da capacidade local de implementação e das características produtivas, biofísicas e geográficas do território. A \textbf{intertemporalidade} integra a motivação do estudo porque as decisões presentes de conversão respondem a retornos econômicos esperados e modificam o estoque de vegetação disponível nos períodos seguintes, alterando as possibilidades futuras de uso da terra. 

A contribuição pretendida ao debate teórico consiste em avaliar em que medida as relações entre retornos econômicos esperados, instituições e conversão da vegetação nativa identificadas pela literatura se sustentam em diferentes condições territoriais. Mato Grosso oferece um caso especialmente informativo para esse exame por reunir produção agropecuária em larga escala e três biomas -- Amazônia, Cerrado e Pantanal -- com distintas características biofísicas e trajetórias de uso da terra. Ao comparar atividades agropecuárias e biomas sob uma base territorial comum, o artigo busca esclarecer o alcance dessas relações e as condições que modificam sua intensidade. A análise da interação entre preços, capacidade produtiva local e ambiente institucional poderá, assim, ajudar a compreender por que incentivos econômicos semelhantes se associam a respostas distintas de conversão, oferecendo evidências para discutir os limites de aplicação das explicações existentes.

%% Relevância aplicada e contribuição para políticas: decisões privadas, fiscalização pública, governança e limites institucionais.
Pragmaticamente, compreender essas relações pode oferecer subsídios para: (a) decisões de firmas e produtores sobre usos alternativos da terra, considerando os incentivos econômicos e as restrições institucionais associados à conversão e à conservação; (b) a discussão de mecanismos de mercado que busquem compatibilizar produção agropecuária e conservação da vegetação nativa, sem constituir uma avaliação direta da eficácia desses mecanismos; e (c) ações públicas de desenho de políticas, fiscalização e iniciativas de conservação e de alocação de recursos ambientais por governos e organizações do terceiro setor.

Isto posto, o artigo busca contribuir ao distinguir a conversão bruta da vegetação nativa para agricultura e pastagem da reorganização produtiva em áreas anteriormente convertidas. Essa distinção permite examinar em quais condições os incentivos econômicos se associam à abertura de novas áreas e como essa resposta varia com a capacidade produtiva local e o ambiente institucional. Articulada a H1, a análise poderá esclarecer como as decisões de uso da terra modificam o estoque de vegetação nativa, oferecendo evidência sobre sua formação como componente endógeno do processo econômico e subsídios à discussão de políticas de conservação que considerem a heterogeneidade territorial.

% ------------------------------------
%%% (IV) METODOLOGIA
% ------------------------------------

%% Estratégia empírica: dados em painel (2001–2024), fontes, ordenação temporal e identificação econométrica.
H1 será examinada com dados em painel para Mato Grosso, tendo como candidatos a unidade de análise município-ano e o período 2001-2024. A análise combinará transições anuais de uso da terra, preços agropecuários, medidas predeterminadas de capacidade produtiva e registros institucionais. As especificações considerarão a ordenação temporal entre incentivos e conversão, efeitos fixos territoriais e temporais e controles compatíveis com a estratégia de identificação. Esses procedimentos não asseguram, por si sós, identificação causal. Quando não for possível sustentar as condições necessárias a essa identificação, os resultados serão interpretados como associações condicionais. A endogeneidade da fiscalização, a dependência espacial e a validade temporal das medidas de capacidade produtiva serão consideradas na definição das especificações e na interpretação dos resultados. As comparações entre atividades e biomas dependerão do suporte amostral disponível. A definição da amostra final depende da harmonização territorial e da obtenção das conversões brutas, ainda pendentes.

% DICA: Usar: GoldsmithPinkhamEtAl2020

% ------------------------------------
%%% (V) ORGANIZAÇÃO DO ARTIGO
% ------------------------------------

%% Organização do artigo: estrutura das seções subsequentes (1 a 5).
Além desta introdução, o artigo está organizado em cinco seções. A primeira revisa a literatura sobre retornos, instituições e conversão da vegetação. A segunda apresenta Mato Grosso como contexto empírico, descreve os dados e explicita a construção das variáveis. A terceira desenvolve a estratégia econométrica para examinar H1. A quarta estabelece os critérios de apresentação e discussão dos resultados a serem obtidos, com atenção à heterogeneidade e à robustez. Por fim, a quinta reúne as contribuições pretendidas e as limitações do estudo. Nesta versão para qualificação, não são apresentados resultados econométricos.

Este artigo concentra-se nos retornos esperados dos usos agropecuários e em seu condicionamento institucional, ficando a análise dos mecanismos de apropriação irregular e valorização especulativa da terra reservada a artigo futuro.